\documentclass{article}
\input{../../../lib.sty}

\title{\texttt{fn make\_stable\_group\_by}}
\author{\MichaelShoemateAuthor}
\begin{document}
\maketitle


\contrib

Proves soundness of \rustdoc{transformations/make_stable_lazyframe/group_by/fn}{make\_stable\_group\_by}
in \asOfCommit{mod.rs}{a105afe}.

\section{Hoare Triple}
\subsection*{Precondition}
\subsubsection*{Compiler-verified}
\begin{itemize}
    \item Generic \texttt{M} implements \rustdoc{transformations/cast_metric/trait}{UnboundedMetric}.
\end{itemize}

\subsubsection*{Caller-verified}
None

\subsection*{Pseudocode}
\label{sec:python-pseudocode}
\lstinputlisting[language=Python,firstline=2,escapechar=|]{./pseudocode/make_stable_group_by.py}

\subsection*{Postcondition}
\begin{theorem}
\validTransformation{\texttt{(input\_domain, input\_metric, plan)}}{\\\texttt{make\_stable\_group\_by}}
\end{theorem}

\begin{proof}[Appropriate Output Domain]
    By line \ref{line:stable-keys} the grouping keys are stable row-by-row transformations of the data.
    By line \ref{line:infallible-aggs} the aggregates are infallible.
    This validation also rejects physical row indices and positional Series/range
    literals, and requires a public, finite scalar quantile parameter in $[0,1]$:
    Polars evaluates that parameter
    against the whole frame, so data-dependent quantiles would not be group-local.
    Therefore the accepted aggregates do not raise data-dependent errors or
    access those non-local inputs.

    By the postcondition of \texttt{DslPlan.simulate\_schema},
    \texttt{series\_domains} follows the expected schema of members in the output domain.
    The output column domains preserve the schema while allowing infallible black-box aggregation.

    Only margins whose original grouping expressions are a subset of the grouping keys are preserved.
    On line \ref{line:materialized-margin}, each preserved key becomes \texttt{col(output\_name)}.
    Polars evaluates a grouping expression on the input and stores its value in the output column with that name.
    The translated margin therefore partitions output rows by the same labels as the original margin partitions input rows.
    Each input row contributes to at most one output row,
    so the upper bounds on group lengths and the number of groups remain valid.
    Re-evaluating the original expression after aggregation is not generally equivalent
    and must not inherit the original margin.
    Among these margins, invariants are discarded.

    It has been shown that for every element $x$ in \texttt{input\_domain}, $\function(x)$ is in \texttt{output\_domain} or raises a data-independent runtime exception.
\end{proof}


The stability guarantee doesn't attempt to claim the broadest set of possible bounds on output distances,
rather it only claims a simple bound that might be useful for the user.
This can be extended in the future, but is sufficient for select queries.

\begin{proof}[Stability guarantee]
    We first simplify the problem on line \ref{line:identifier} by only considering datasets that differ by rows, not by identifiers.

    The output metric is always \texttt{FrameDistance<SymmetricDistance>},
    regardless of the input metric or \texttt{maintain\_order} option.
    Thus arbitrary output-group reorderings do not contribute distance.
    The only output distance bound derived is for the empty grouping set, namely the entire output table.

    For unordered inputs, zero-distance permutations must preserve the joint
    distribution of the aggregates. Line \ref{line:permutation-invariance}
    enforces a conservative expression allowlist for this purpose.
    Scalar literals broadcast without depending on row positions; positional
    Series and range literals do not satisfy that property.
    Accepted rowwise expressions commute with row permutations.
    Grouping keys must also have canonical equality representatives; unordered
    grouping on float-containing keys is rejected because Polars may emit either
    signed-zero representation. Invariant reducers, such as integer sums, therefore
    produce the same values. Floating-point min/max are not treated as invariant:
    Polars' sorted-column fast paths may select bit-distinct tied values.
    Sensitive reducers, including floating-point sums and means, instead require
    a trusted uniform random permutation of their input. A generic Polars sort
    is not treated as canonical: the engine may skip sorting based on existing
    sorted-state metadata that this analysis does not certify.
    For uniform permutations, couple corresponding records in the two inputs;
    the resulting ordered samples agree in distribution, including when only a
    prefix is retained. Deterministic postprocessing preserves this distributional
    equality. Independently randomized expressions can be coupled independently.
    For ordered inputs, unchanged groups already have the same relative input order.
    In either case, unchanged groups can be coupled to have identical aggregates.

    We first retrieve optional upper bounds on the number of rows and groups an individual may contribute on lines \ref{line:total-contributions} and \ref{line:total-groups}.
    Both of these bounds directly correspond to the number of rows an individual may influence in the resulting dataset.

    \begin{align}
        &\max_{x \sim x'} d_{\mathrm{FrameDistance<SymmetricDistance>}}(\texttt{function}(x), \texttt{function}(x')) \\
        &\le \mathrm{option\_min}(\texttt{contributed\_rows}, \texttt{contributed\_groups}) \cdot 2
    \end{align}

    Since each influenced row accounts for one addition and one removal, the distance is twice the number of influenced rows.
    This is reflected on line \ref{line:double}.
    The resulting bound is returned on line \ref{line:bound}, satisfying the postcondition.
\end{proof}
\end{document}
